\abstract{%
Geometry-first approaches to spacetime reconstruction aim to explain how macroscopic notions of space, time, and causality can emerge from mathematically controlled structures without presupposing a continuum background time parameter. This review synthesizes a family of constructions organized around a concrete dependency chain: (i) six-dimensional superspace and cut-and-project model sets as geometric primitives for quasi-periodic ``screens''; (ii) finite-resolution \emph{readout} modeled as a windowed projection protocol that maps continuous internal degrees of freedom to symbolic data streams; (iii) stabilization/canonicalization maps (Fold-type) that compress raw readout words into robust constrained languages and thereby induce preimage fibers encoding unresolved degrees of freedom; and (iv) an information-theoretic notion of time defined by entropy-rate-controlled refinement of posterior cylinder measures (Shannon--McMillan--Breiman / Kolmogorov--Sinai framework). We highlight how standard superspace concepts (acceptance domains, atomic surfaces) interact with symbolic codings (mechanical/Sturmian processes) to provide an auditable ``geometry-to-data'' interface. A central theme is the strict separation between object-level geometry, protocol-level definitions, and physical interpretation: claims are consistently labeled as definitions, assumptions, or derived propositions. Within this separation we formulate a conservative causality interface---including a lightcone-boundary readout viewpoint---and clarify what would be required to relate protocol precedence to Lorentzian causality. We conclude with falsifiable fingerprints (quasiperiodic spectral modules, resolution-dependent visibility curves, and folding/degeneracy counts) and with open problems for bridging superspace readout models to classical and quantum spacetime geometry.%
}

\keywords{cut-and-project; superspace; model sets; symbolic dynamics; entropy rate; causality}

